\documentclass[10pt,a4paper]{amsart}
\usepackage[margin=19mm]{geometry}
\usepackage{amsmath,amssymb,mathtools}
\usepackage[scr=rsfs,cal=euler]{mathalfa}
\usepackage{xfrac}
\usepackage[unicode,hidelinks,bookmarks=false]{hyperref}
\newcommand{\ie}{i.e.\ }
\newcommand{\cf}{cf.\ }
\newcommand{\resp}{respectively\ }
\newcommand{\Subsection}{\subsection}
\providecommand{\nobreakdash}{\nobreak\hbox{-}}
\setlength{\emergencystretch}{2em}
\multlinegap0pt
\usepackage{amssymb}
\usepackage{xcolor}
\newtheorem{lemma}{Lemma}
\newcommand{\m}{\mbox{\textnormal{mid}}}
\title{Long-time asymptotics of (1,3)-sign solitary waves for the damped nonlinear Klein-Gordon equation}
\author{Kenjiro Ishizuka}
\date{Source: arXiv:2602.01205v2 (CC BY 4.0)}
\begin{document}
\maketitle

\noindent\emph{Source and licence.} Long-time asymptotics of (1,3)-sign solitary waves for the damped nonlinear Klein-Gordon equation. Version arXiv:2602.01205v2 (CC BY 4.0), distributed by arXiv under the Creative Commons Attribution 4.0 International licence, http://creativecommons.org/licenses/by/4.0/, as recorded in the arXiv licence metadata for this exact version and shown on the abstract page https://arxiv.org/abs/2602.01205v2. Full licence notice: \texttt{licenses/arxiv-2602.01205-CC-BY-4.0.txt}.\par\smallskip

\noindent\textit{Editorial scope.} Counted unit: the article's lemma minhenchou (source0.tex lines 1653--1662). The article's complete original proof of it is delivered with it (source0.tex lines 1663--1804, one \texttt{proof} environment of 142 lines). No mathematics is written, repaired or re-derived: the statement and the proof are the article's own lines, in the article's own notation, and no retained line is substituted or re-flowed.  Retained uncounted as the source-local context the proof needs: lines 529--531; lines 667--669; lines 1095--1128; lines 1540--1549.  Nothing was omitted from any retained span, so the package carries no omission record.  No retained line contains a citation, so the wrapper carries no bibliography and no bibliography entry.  No retained line contains a figure environment, an image-inclusion command or drawing code, so no image is delivered and the package declares no asset dependency.  Editorial formatting only, in addition: an AMS \texttt{amsart} preamble that restates the article's own declarations and macros used by the retained lines, namely the environments \texttt{lemma} on separate counters, together with the standard packages those lines need.  The retained environments print in this document as Lemma 1 in order of appearance, whereas the article prints the same items with the numbering of its own full document.  The complete native source of the article as distributed in the arXiv e-print for this exact version is bundled unchanged as \texttt{sources/193774/source0.tex}.
\begin{align}\label{tukaeruF1}
\left|\mathcal{F}(R+a)-e^{-a}\mathcal{F}(R)\right|\lesssim \frac{\mathcal{F}(R)}{R},
\end{align}

\begin{align}\label{Apositive}
A(t)\geq 0.
\end{align}

\begin{lemma}\label{sankakukeilemma}
We assume $M\gg 1$. Then, $\varsigma_1,\varsigma_2,\varsigma_3\in\mathbb{R}^d$ satisfy the following.
\begin{enumerate}
\item When $\varsigma_1,\varsigma_2,\varsigma_3$ satisfy
\begin{align*}
\left|\ |\varsigma_1-\varsigma_2|-M\ \right|&\leq M^{\frac{99}{100}},
\\
\left|\ |\varsigma_1-\varsigma_3|-M\ \right|&\leq M^{\frac{99}{100}},
\\
|\varsigma_2-\varsigma_3|-M&\geq  -M^{\frac{99}{100}},
\end{align*}
we have
\begin{align*}
\frac{\varsigma_2-\varsigma_1}{|\varsigma_2-\varsigma_1|}\cdot \frac{\varsigma_3-\varsigma_1}{|\varsigma_3-\varsigma_1|}-\frac{1}{2}&\leq 5M^{-\frac{1}{100}},
\\
\frac{\varsigma_1-\varsigma_2}{|\varsigma_1-\varsigma_2|}\cdot \frac{\varsigma_3-\varsigma_2}{|\varsigma_3-\varsigma_2|}-\frac{1}{2}&\geq -5M^{-\frac{1}{100}},
\\
\frac{\varsigma_1-\varsigma_3}{|\varsigma_1-\varsigma_3|}\cdot \frac{\varsigma_2-\varsigma_3}{|\varsigma_2-\varsigma_3|}-\frac{1}{2}&\geq -5M^{-\frac{1}{100}}.
\end{align*}

\item When $\varsigma_1,\varsigma_2,\varsigma_3$ satisfy
\begin{align*}
\left|\ |\varsigma_1-\varsigma_2|-M\ \right|&\leq M^{\frac{99}{100}},
\\
|\varsigma_1-\varsigma_3|-M&\geq -M^{\frac{99}{100}},
\\
|\varsigma_2-\varsigma_3|-M&\geq -M^{\frac{99}{100}},
\end{align*}
we have
\begin{align*}
\frac{\varsigma_1-\varsigma_3}{|\varsigma_1-\varsigma_3|}\cdot \frac{\varsigma_2-\varsigma_3}{|\varsigma_2-\varsigma_3|}-\frac{1}{2}&\geq -5M^{-\frac{1}{100}}.
\end{align*}
\end{enumerate}
\end{lemma}

\begin{lemma}\label{midhenchou}
Let $\vec{u}$ be a (1,3)-sign solitary waves. Furthermore, we assume that we have for $T_1\leq t\leq T_2$
\begin{align}\label{estimate0125}
\tilde{D}(t)-\hat{D}(t)<D^{\frac{1}{4}}(t).
\end{align}
Then, we have for $T_1\leq t\leq T_2$
\begin{align}\label{DA21}
\m{(\rho_{12}(t),\rho_{13}(t),\rho_{23}(t))}\geq \tilde{D}(t)+D^{\frac{1}{3}}(t).
\end{align}
\end{lemma}

\begin{lemma}\label{minhenchou}
Let $\vec{u}$ be a (1,3)-sign solitary waves . Furthermore, we assume that there exist $1\ll T_1< T_2\leq \infty$ such that we have for $T_1\leq t<T_2$
\begin{align}
\tilde{D}(t)-\hat{D}(t)<D^{\frac{1}{4}}(t). \label{lemass01253}
\end{align}
Then, we have for $T_1\leq t<T_2$
\begin{align}\label{Dmodes1}
\left(D_{mod}(t)-\tilde{D}(t)\right)-\left(D_{mod}(T_1)-\tilde{D}(T_1)\right)\gtrsim \int_{T_1}^t \mathcal{F}(D(s))ds.
\end{align}
\end{lemma}

\begin{proof}
By symmetry, we may assume that $\tilde{D}(T_1)=\rho_{12}(T_1)$.
In this case, by Lemma \ref{midhenchou} and continuity, we have for $T_1\leq t<T_2$.
\begin{align*}
\rho_{12}(t)=\tilde{D}(t),\ \rho_{13}(t)>\tilde{D}(t)+D^{\frac{1}{3}}(t),\ \rho_{23}(t)>\tilde{D}(t)+D^{\frac{1}{3}}(t).
\end{align*}
Then we have
\begin{align*}
\dot{\rho_1}&=2\mathcal{F}(\rho_1)+\mathcal{F}(\rho_2)c_{12}+\mathcal{F}(\rho_3)c_{13}+\mathcal{F}(\rho_{12})u_{12}\cdot u_1+o\left(\mathcal{F}(D)\right),
\\
\dot{\rho}_2&=\mathcal{F}(\rho_1)c_{12}+2\mathcal{F}(\rho_2)+\mathcal{F}(\rho_3)c_{23}+\mathcal{F}(\rho_{12})u_{21}\cdot u_2+o\left(\mathcal{F}(D)\right),
\\
\dot{\rho}_3&=\mathcal{F}(\rho_{1})c_{13}+\mathcal{F}(\rho_2)c_{23}+2\mathcal{F}(\rho_3)+o\left(\mathcal{F}(D)\right),
\\
\dot{\rho}_{12}&=-2\mathcal{F}(\rho_{12})+\mathcal{F}(\rho_2)(u_{12}\cdot u_2)+\mathcal{F}(\rho_1)(u_1\cdot u_{21})+o\left(\mathcal{F}(D)\right).
\end{align*}
We note that by definition we have $|D_{mod}-\hat{D}|\leq c_{\heartsuit}$.
First, we assume that $D_{mod}=\rho_1$.
Since we have $\rho_1\leq \rho_3+c_{\heartsuit}$, by \eqref{tukaeruF1} we obtain
\begin{align}\label{rho1rho3hyouka}
\mathcal{F}(\rho_3)-e^{c_{\heartsuit}}\mathcal{F}(\rho_1)\lesssim \frac{\mathcal{F}(D)}{D}.
\end{align}
Since we have for $T_1\leq t<T_2$
\begin{align*}
|\rho_1-D|+|\rho_{12}-D|\leq D^{\frac{99}{100}},
\\
\rho_{2}-D\geq -D^{\frac{99}{100}},
\end{align*}
by Lemma \ref{sankakukeilemma} we obtain
\begin{align}\label{kakudo1261}
c_{12}-\frac{1}{2}\gtrsim -D^{-\frac{1}{100}},\ u_1\cdot u_{21}-\frac{1}{2}\lesssim D^{-\frac{1}{100}}.
\end{align}
Furthermore, if we have $\rho_2>D+D^{\frac{99}{100}}$, then we have
\begin{align*}
\mathcal{F}(\rho_2)\lesssim  \frac{\mathcal{F}(D)}{D},
\end{align*}
and therefore we obtain
\begin{align*}
\dot{\rho}_1-\dot{\rho}_{12}&\geq \frac{3}{2}\mathcal{F}(\rho_1)-\mathcal{F}(\rho_3)+o\left(\mathcal{F}(D)\right).
\end{align*}
On the other hand, if we have $\rho_2<D+D^{\frac{99}{100}}$, then by Lemma \ref{sankakukeilemma} we have
\begin{align*}
\left|c_{12}-\frac{1}{2}\right|+\left|u_{12}\cdot u_2-\frac{1}{2}\right|+\left|u_1\cdot u_{21}-\frac{1}{2}\right|\lesssim D^{-\frac{1}{100}},
\end{align*}
and therefore we obtain
\begin{align*}
\dot{\rho}_1-\dot{\rho}_{12}&\geq \frac{3}{2}\mathcal{F}(\rho_1)-\mathcal{F}(\rho_3)+o\left(\mathcal{F}(D)\right).
\end{align*}
Hence, in either case,
\begin{align}\label{rho1-rho12}
\dot{\rho}_1-\dot{\rho}_{12}&\geq \frac{3}{2}\mathcal{F}(\rho_1)-\mathcal{F}(\rho_3)+o\left(\mathcal{F}(D)\right).
\end{align}
By \eqref{kakudo1261} and \eqref{rho1-rho12}, we have
\begin{align}\label{rho1minnnotoki}
\begin{aligned}
\dot{\rho}_1-\dot{\rho}_{12}&\geq\frac{3}{2}\mathcal{F}(\rho_1)-\frac{4+\sqrt{3}}{4}\mathcal{F}(\rho_1)+o\left(\mathcal{F}(D)\right)
\\
&=\frac{2-\sqrt{3}}{4}\mathcal{F}(D_{mod})+o\left(\mathcal{F}(D)\right)
\\
&\gtrsim \mathcal{F}(D).
\end{aligned}
\end{align}
The case $D_{mod}=\rho_2$ can be treated by the same argument as above. Next, we assume that $D_{mod}=\rho_3+c_{\heartsuit}$. When $\rho_1$ and $\rho_2$ satisfy
\begin{align*}
\rho_1&>D+D^{\frac{99}{100}},
\\
\rho_2&>D+D^{\frac{99}{100}},
\end{align*}
we have
\begin{align*}
\dot{\rho}_3-\dot{\rho}_{12}\geq 2\mathcal{F}(\rho_3)+o\left(\mathcal{F}(D)\right)\gtrsim \mathcal{F}(D).
\end{align*}
Next, if $\rho_1$ and $\rho_2$ satisfy 
\begin{align*}
\rho_1&<D+D^{\frac{99}{100}},
\\
\rho_2&>D+D^{\frac{99}{100}},
\end{align*}
by Lemma \ref{sankakukeilemma}, we have
\begin{align*}
u_1\cdot u_{21}-\frac{1}{2}\lesssim D^{-\frac{1}{100}}.
\end{align*}
Therefore we obtain
\begin{align*}
\dot{\rho}_3-\dot{\rho}_{12}&\geq 2\mathcal{F}(\rho_3)-\frac{3}{2}\mathcal{F}(\min{(\rho_1,\rho_2)})+o\left(\mathcal{F}(D)\right)\gtrsim \mathcal{F}(D).
\end{align*}
When $\rho_1, \rho_2$ satisfy $\rho_1>D+D^{\frac{99}{100}}$ and $\rho_2<D+D^{\frac{99}{100}}$, the same argument yields the above inequality.

Last, when $\rho_1$ and $\rho_2$ satisfy
\begin{align*}
\rho_1&<D+D^{\frac{99}{100}},
\\
\rho_2&<D+D^{\frac{99}{100}},
\end{align*}
by Lemma \ref{sankakukeilemma} we have
\begin{align}\label{DA31}
\left|c_{12}-\frac{1}{2}\right|+\left|u_{12}\cdot u_2-\frac{1}{2}\right|+\left|u_1\cdot u_{21}-\frac{1}{2}\right|\lesssim D^{-\frac{1}{100}}.
\end{align}
By \eqref{Apositive} and \eqref{DA31}, we have
\begin{align*}
\frac{3}{4}+c_{13}c_{23}-c_{13}^2-c_{23}^2\gtrsim -D^{-\frac{1}{100}}.
\end{align*}
Furthermore, since we have
\begin{align*}
c_{13}^2+c_{23}^2-c_{13}c_{23}=\frac{1}{4}\left(c_{13}+c_{23}\right)^2+\frac{3}{4}\left(c_{13}-c_{23}\right)^2\geq \frac{1}{4}\left(c_{13}+c_{23}\right)^2,
\end{align*}
we obtain
\begin{align}\label{c13c23es1}
c_{13}+c_{23}+\sqrt{3}\gtrsim -D^{-\frac{1}{100}}.
\end{align}
We now combine the arguments above. First, if $c_{13}, c_{23}$ satisfy $c_{13}\geq 0$ or $c_{23}\geq 0$, then we have
\begin{align*}
\dot{\rho}_3-\dot{\rho}_{12}&\geq 2\mathcal{F}(\rho_3)-2\mathcal{F}(\min{(\rho_1,\rho_2)})+o\left(\mathcal{F}(D)\right)
\\
&\gtrsim \mathcal{F}(D).
\end{align*}
If $c_{13}, c_{23}$ satisfy $c_{13}<0$ and $c_{23}<0$, then by \eqref{c13c23es1} we have
\begin{align}\label{rho3-rho12}
\begin{aligned}
\dot{\rho}_3-\dot{\rho}_{12}&\geq 2\mathcal{F}(\rho_3)+(c_{13}+c_{23}-1)\mathcal{F}(\min{(\rho_1,\rho_2)})+o\left(\mathcal{F}(D)\right)
\\
&\geq 2\mathcal{F}(\rho_3)-(\sqrt{3}+1)\mathcal{F}(\min{(\rho_1,\rho_2)})+o\left(\mathcal{F}(D)\right).
\end{aligned}
\end{align}
Furthermore, since $\rho_3+c_{\heartsuit}\leq \min{(\rho_1,\rho_2)}$ holds, we obtain
\begin{align}\label{minrho1rho2es}
\mathcal{F}(\min{(\rho_1,\rho_2)})\leq \frac{4}{4+\sqrt{3}}\mathcal{F}(\rho_3)+o\left(\mathcal{F}(D)\right).
\end{align}
By \eqref{rho3-rho12} and \eqref{minrho1rho2es}, we have
\begin{align*}
\dot{\rho}_3-\dot{\rho}_{12}&\geq \left(2-\frac{4(\sqrt{3}+1)}{4+\sqrt{3}}\right)\mathcal{F}(\rho_3)+o\left(\mathcal{F}(D)\right)
\\
&=\frac{4-2\sqrt{3}}{4+\sqrt{3}}\mathcal{F}(\rho_3)+o\left(\mathcal{F}(D)\right)
\\
&\gtrsim \mathcal{F}(D).
\end{align*}
From the above, we have for $T_1\leq t<T_2$,
\begin{align}\label{Dmodes}
\frac{d}{dt}D_{mod}-\dot{\rho}_{12}\gtrsim \mathcal{F}(D).
\end{align}
Hence, integrating \eqref{Dmodes} on $[T_1,t]$ yields \eqref{Dmodes1} and completes the proof.
\end{proof}

\end{document}
